% !TeX program = xelatex
% !TeX encoding = UTF-8
% Processed from chunks/02_lindberg_beginnings.md
\documentclass[UTF8]{ctexart}
\usepackage[a4paper,margin=2.5cm]{geometry}
\usepackage{graphicx}
\usepackage{amsmath}
\usepackage{caption}
\setlength{\parindent}{0pt}
\setlength{\parskip}{0.75em}

\newcommand{\orig}{\par\noindent\textbf{Original.}\quad}
\newcommand{\zh}{\par\noindent\textbf{译文。}\quad}

\begin{document}

\section*{The Beginnings of Western Science / 《西方科学的起源》}
\textbf{David C. Lindberg / 戴维 C. 林德伯格}

\section*{Chapter Two: The Greeks and the Cosmos / 第二章：希腊人与宇宙}

\subsection*{Plato's World of Forms / 柏拉图的形式世界}

\orig The death of Socrates in 399 B.C., coming as it did around the turn of the century (not on their calendar, of course, but on ours), has made it a convenient point of demarcation in the history of Greek philosophy. Thus Socrates' predecessors of the sixth and fifth centuries (the philosophers who have occupied us until now in this chapter) are commonly called the ``pre-Socratic philosophers.'' But Socrates' prominence is more than an accident of the calendar, for Socrates represents a shift in emphasis within Greek philosophy, away from the cosmological concerns of the sixth and fifth centuries toward political and ethical matters. Nonetheless, the shift was not so dramatic as to preclude continuing attention to the major problems of pre-Socratic philosophy. We find both the new and the old in the work of Socrates' younger friend and disciple, Plato (fig. 2.4).

\zh 苏格拉底死于公元前399年，恰好接近一个世纪之交（当然不是按他们的历法，而是按我们的历法），因此成了划分希腊哲学史的一个方便界标。于是，公元前六、五世纪那些早于苏格拉底的哲学家（也就是本章迄今讨论的那些哲学家）通常被称为“前苏格拉底哲学家”。但是，苏格拉底的重要性并不只是历法上的偶然；他代表了希腊哲学重心的一次转移：从公元前六、五世纪的宇宙论关怀，转向政治和伦理问题。尽管如此，这种转移并没有剧烈到使人们不再关注前苏格拉底哲学的主要问题。在苏格拉底较年轻的朋友和弟子柏拉图的著作中，我们既能看到新的关怀，也能看到旧的问题（图2.4）。

\begin{figure}[htbp]
\centering
\includegraphics[width=0.42\linewidth]{../chunks/images/76f2553df76f729a90b5a54f134b9c18710e62ab8ac46e3c53b24a76fd0f8cbb.jpg}
\caption{Fig. 2.4. Plato (1st c. A.D. copy). Museo Vaticano, Vatican City. Alinari/Art Resource N.Y. / 图2.4 柏拉图（公元1世纪摹本）。梵蒂冈城，梵蒂冈博物馆。Alinari/Art Resource N.Y.}
\end{figure}

\orig Plato (427--348/47 B.C.) was born into a distinguished Athenian family, active in affairs of state; he was undoubtedly a close observer of the political events that led up to Socrates' execution. After Socrates' death, Plato left Athens and visited Italy and Sicily, where he seems to have come into contact with Pythagorean philosophers. In 388 Plato returned to Athens and founded a school of his own, the Academy, where young men could pursue advanced studies (see fig. 4.1). Plato's literary output appears to have consisted almost entirely of dialogues, the majority of which have survived. We will find it necessary to be highly selective in our examination of Plato's philosophy; let us begin with his quest for the underlying reality.

\zh 柏拉图（公元前427--348/47年）出生于一个显赫的雅典家族，这个家族积极参与城邦政治；毫无疑问，他密切目睹了导致苏格拉底被处死的一系列政治事件。苏格拉底死后，柏拉图离开雅典，游历意大利和西西里；在那里，他似乎接触到了毕达哥拉斯学派的哲学家。公元前388年，柏拉图回到雅典，创办了自己的学校“学园”（Academy），年轻人在那里可以从事高等研究（见图4.1）。柏拉图的著作似乎几乎全由对话构成，其中大部分保存至今。我们在考察柏拉图哲学时必须高度取舍；先从他对根本实在的追寻说起。

\begin{figure}[htbp]
\centering
\includegraphics[width=0.70\linewidth]{../chunks/images/5dec5d9ba0cfe21df1e401825ccfb025e409a982dea1eaad5bbd1158e8374444.jpg}
\caption{Fig. 4.1. The schools of Hellenistic Athens. \copyright{} Candace H. Smith. / 图4.1 希腊化时期雅典的各学派。\copyright{} Candace H. Smith.}
\end{figure}

\orig In a passage in one of his dialogues, the \textit{Republic}, Plato reflected on the relationship between the actual tables constructed by a carpenter and the idea or definition of a table in the carpenter's mind. The carpenter replicates the mental idea as closely as possible in each table he makes, but always imperfectly. No two manufactured tables are alike down to the smallest detail, and limitations in the material (a knot here, a warped board there) ensure that none will fully measure up to the ideal.

\zh 在其对话《理想国》的一段文字中，柏拉图思考了木匠实际制作的桌子与木匠心中“桌子”的理念或定义之间的关系。木匠在每一张桌子上都尽可能复制心中的理念，但这种复制总是不完美。没有两张制成的桌子在最细微之处完全相同，材料的局限（这里有个木节，那里有块弯曲的木板）也保证了没有哪张桌子能完全达到理想。

\orig Now, Plato argued, there is a divine craftsman who bears the same relationship to the cosmos as the carpenter bears to his tables. The divine craftsman (the Demiurge) constructed the cosmos according to an idea or plan, so that the cosmos and everything in it are replicas of eternal ideas or forms---but always imperfect replicas because of limitations inherent in the materials available to the Demiurge. In short, there are two realms: a realm of forms or ideas, containing the perfect form of everything; and the material realm in which these forms or ideas are imperfectly replicated.

\zh 柏拉图于是主张，有一位神圣工匠与宇宙的关系，正如木匠与其桌子的关系一样。这位神圣工匠（德谟革，Demiurge）按照某种理念或计划构造了宇宙；因此，宇宙以及其中的一切都是永恒理念或形式的摹本。但由于德谟革所能使用的材料本身有限，这些摹本总是不完美。简言之，存在两个领域：一个是形式或理念的领域，包含万物的完美形式；另一个是质料领域，在其中，这些形式或理念被不完美地复制出来。

\orig Plato's notion of two distinct realms will seem strange to many people, and we must therefore stress several points of importance. The forms are incorporeal, intangible, and insensible; they have always existed, sharing the property of eternity with the Demiurge; and they are absolutely changeless. They include the form, the perfect idea, of everything in the material world. One does not speak of their location, since they are incorporeal and therefore not spatial. Although incorporeal and imperceptible by the senses, they objectively exist; indeed, true reality (reality in its fullness) is located only in the world of forms. The sensible, corporeal world, by contrast, is imperfect and transitory. It is less real in the sense that the corporeal object is a replica of, and therefore dependent for its existence upon, the form. The form has primary existence, its corporeal replica secondary existence.

\zh 柏拉图关于两个不同领域的观念，对许多人来说会显得陌生，因此我们必须强调几个要点。形式是无形体、不可触摸、不可感知的；它们始终存在，与德谟革一样具有永恒性；并且绝对不变。它们包括质料世界中一切事物的形式，也就是其完美理念。我们不能谈论它们的位置，因为它们没有形体，因而也不占据空间。尽管形式无形体且不能由感官感知，它们却客观存在；事实上，真正的实在（充分意义上的实在）只位于形式世界。相比之下，可感的、有形的世界是不完美且短暂的。它之所以较少真实，是因为有形事物只是形式的摹本，因此其存在依赖于形式。形式具有首要存在，它的有形摹本则具有次级存在。

\orig Plato illustrated this conception of reality in his famous ``allegory of the cave,'' found in book VII of the \textit{Republic}. Men are imprisoned within a deep cave, chained so as to be incapable of moving their heads. Behind them is a wall, and beyond that a fire. People walk back and forth behind the wall, holding above it various objects, including statues of humans and animals; the objects cast shadows on the wall that is visible to the prisoners. The prisoners see only the shadows cast by these objects; and, having lived in the cave from childhood, they no longer recall any other reality. They do not suspect that these shadows are but imperfect images of objects that they cannot see; and consequently they mistake the shadows for the real.

\zh 柏拉图在《理想国》第七卷著名的“洞穴寓言”中说明了这种实在观。人们被囚禁在一个深洞之中，身上缚有锁链，头也无法转动。他们身后有一堵墙，墙后有一团火。有人在墙后往来行走，高举着各种物体，包括人和动物的雕像；这些物体把影子投到囚徒能够看见的墙面上。囚徒只能看见这些物体投下的影子；由于他们自幼便生活在洞中，早已不记得任何别的实在。他们并不怀疑这些影子只是他们看不见的物体的不完美影像；于是，他们误把影子当成真实。

\orig So it is with all of us, says Plato. We are souls imprisoned in bodies. The shadows of the allegory represent the world of sense experience. The soul, peering out from its prison, is able to perceive only these flickering shadows, and the ignorant claim that this is all there is to reality. However, there do exist the statues and other objects of which the shadows are feeble representations and also the humans and animals of which the statues are imperfect replicas. To gain access to these higher realities, we must escape the bondage of sense experience and climb out of the cave, until we find ourselves able, finally, to gaze on the eternal realities, thereby entering the realm of true knowledge.

\zh 柏拉图说，我们所有人都是如此。我们是被囚禁在身体中的灵魂。寓言中的影子代表感官经验的世界。灵魂从它的牢笼中向外窥视，只能感知这些闪烁不定的影子；无知者便声称这就是实在的全部。然而，确实存在着那些雕像和其他物体，影子只是它们微弱的表象；也确实存在着人和动物，雕像只是它们不完美的摹本。若要接近这些更高的实在，我们必须摆脱感官经验的束缚，爬出洞穴，直到最终能够凝视那些永恒实在，从而进入真正知识的领域。

\orig What are the implications of these views for the concerns of the pre-Socratic philosophers? First, Plato equated his forms with the underlying reality, while assigning derivative or secondary existence to the corporeal world of sensible things. Second, Plato has made room for both change and stability by assigning each to a different level of reality: the corporeal realm is the scene of imperfection and change, while the realm of forms is characterized by eternal, changeless perfection. Both change and stability are therefore genuine; each characterizes something; but changelessness belongs to the forms and thus shares their fuller reality.

\zh 这些观点对于前苏格拉底哲学家的关切意味着什么？第一，柏拉图把形式等同于根本实在，同时把可感事物的有形世界置于派生的或次级的存在地位。第二，柏拉图通过把变化和稳定分别安置在不同层次的实在中，为二者都留出了位置：有形领域是不完美和变化的场所，而形式领域则以永恒不变的完美为特征。因此，变化和稳定都是真实的；二者各自表征某种事物；但不变性属于形式，因而分享形式那种更充分的实在性。

\orig Third, as we have seen, Plato addressed epistemological questions, placing observation and true knowledge (or understanding) in opposition. Far from leading upward to knowledge or understanding, the senses are chains that tie us down; the route to knowledge is through philosophical reflection. This is explicit in the \textit{Phaedo}, where Plato maintains the uselessness of the senses for the acquisition of truth and points out that when the soul attempts to employ them it is inevitably deceived.

\zh 第三，如我们所见，柏拉图处理了认识论问题，把观察与真正的知识（或理解）置于对立之中。感官远非引导我们上升到知识或理解的途径，而是束缚我们的锁链；通向知识的道路在于哲学反思。这一点在《斐多篇》中说得很明确：柏拉图主张感官对于获得真理毫无用处，并指出灵魂一旦试图运用感官，就不可避免地受到欺骗。

\orig Now the short account of Plato's epistemology frequently ends here; but there are important qualifications that it would be a serious mistake to omit. Plato did not, in fact, dismiss the senses altogether, as Parmenides had done and as the passage from the \textit{Phaedo} might suggest Plato did. Sense experience, in Plato's view, served various useful functions. First, sense experience may provide wholesome recreation. Second, observation of certain sensible objects (especially those with geometrical properties) may serve to direct the soul toward nobler objects in the realm of forms. Plato used this argument as justification for the pursuit of astronomy. Third, Plato argued (in his theory of reminiscence) that sense experience may actually stir the memory and remind the soul of forms that it knew in a prior existence, thus stimulating a process of recollection that will lead to actual knowledge of the forms.

\zh 对柏拉图认识论的简短说明常常到这里就结束了；但还有一些重要限定，若省略就会造成严重误解。事实上，柏拉图并没有像巴门尼德那样完全排斥感官，尽管《斐多篇》的那段文字可能让人以为他如此。按柏拉图的看法，感官经验具有若干有益功能。第一，感官经验可以提供健康的消遣。第二，观察某些可感对象（尤其是具有几何性质的对象）可以引导灵魂转向形式界中更高贵的对象。柏拉图用这一论证为研究天文学辩护。第三，柏拉图在其“回忆说”中主张，感官经验实际上可以激发记忆，使灵魂想起它在前世曾经认识的形式，从而推动一种回忆过程，并最终导致对形式的真实知识。

\orig Finally, although Plato firmly believed that knowledge of the eternal forms (the highest, and perhaps the only true, form of knowledge) is obtainable only through the exercise of reason, the changeable realm of matter is also an acceptable object of study. Such studies serve the purpose of supplying examples of the operation of reason in the cosmos. If this is what interests us (as it sometimes did Plato), the best method of exploring it is surely to observe it. The legitimacy and utility of sense experience are clearly implied in the \textit{Republic}, where Plato acknowledged that a prisoner emerging from the cave first employs his sense of sight to apprehend living creatures, the stars, and finally the most noble of visible (material) things, the sun. But if he aspires to apprehend ``the essential reality,'' he must proceed ``through the discourse of reason unaided by any of the senses.'' Both reason and sense are thus instruments worth having; which one we employ on a particular occasion will depend on the object of study.

\zh 最后，尽管柏拉图坚定地相信，对永恒形式的知识（最高级、或许也是唯一真正的知识形式）只能通过理性的运用来获得，但变化着的质料领域同样可以成为正当的研究对象。这类研究的目的，是为宇宙中理性运作的方式提供例证。如果我们关心的正是这一点（柏拉图有时确实如此），那么探索它的最佳方法无疑就是观察它。《理想国》清楚地暗示了感官经验的正当性和有用性：柏拉图承认，走出洞穴的囚徒首先运用视觉来把握生物、星辰，最后把握可见（质料性）事物中最高贵者，即太阳。但是，如果他想把握“本质的实在”，就必须“通过不借助任何感官的理性论辩”来前进。因此，理性和感官都是值得拥有的工具；在特定场合使用哪一个，取决于研究对象。

\orig There is another way of expressing all of this, which may shed light on Plato's achievement. When Plato assigned reality to the forms, he was, in fact, identifying reality with the properties that classes of things have in common. The bearer of true reality is not (for example) this dog with the droopy left ear or that one with the menacing bark, but the idealized form of a dog shared (imperfectly, to be sure) by every individual dog---those characteristics by virtue of which we are able to classify all of them as dogs. Therefore, to gain true knowledge, we must set aside all characteristics peculiar to things as individuals and seek the shared characteristics that define them into classes. Now stated in this modest fashion, Plato's view has a distinctly modern ring. Idealization is a prominent feature of a great deal of modern science; we develop models or laws that overlook the incidental in favor of the essential. However, Plato went beyond this, maintaining not merely that true reality is to be found in the common properties of classes of things, but also that this common property (the idea or form) has objective, independent, and indeed prior existence.

\zh 还可以用另一种方式来表述这一切，这或许能显示柏拉图的成就。柏拉图把实在性赋予形式时，实际上是在把实在等同于各类事物共同具有的性质。真正实在的承载者并不是（例如）这条左耳下垂的狗，或那条叫声凶狠的狗，而是所有个别的狗共同分享（诚然，是不完美地分享）的狗的理想化形式，也就是那些使我们能够把它们统统归为狗的特征。因此，若要获得真正知识，我们必须撇开事物作为个体所特有的一切特征，寻求那些把它们界定为某一类别的共同特征。若以这种较为温和的方式表述，柏拉图的观点便颇具现代意味。理想化是许多现代科学的显著特征；我们建立模型或规律时，会忽略偶然因素而偏重本质。不过，柏拉图走得更远：他不仅主张真正实在存在于各类事物的共同性质中，而且主张这种共同性质（理念或形式）具有客观、独立、并且确实先在的存在。

\begin{figure}[htbp]
\centering

\includegraphics[width=0.30\linewidth]{../chunks/images/9519836f4685aaff51f66711e7f2b7597d7605556dc3b35ffcbfe2d575d961d5.jpg}\hfill
\includegraphics[width=0.30\linewidth]{../chunks/images/aa581e862cc9d1691c97016b9b4021f9589afd23337d79179d3b7e7d6cbce561.jpg}\hfill

\includegraphics[width=0.30\linewidth]{../chunks/images/8ea6db286861b81aca334bb0da5448aee2e09ebaf1d510659a5b0cd9f6654773.jpg}
\caption{Uncaptioned image-only material preserved from the OCR source between Chapter 2 and Chapter 3. / OCR 源在第二章与第三章之间保留的无可靠图注图像，按原顺序保留。}
\end{figure}

\section*{Chapter Three: Aristotle's Philosophy of Nature / 第三章：亚里士多德的自然哲学}

\subsection*{Life and Works / 生平与著作}

\orig Aristotle (fig. 3.1) was born in 384 B.C. in the northern Greek town of Stagira, into a privileged family. His father was personal physician to the Macedonian king, Amyntas II (grandfather of Alexander the Great). Aristotle had the advantage of an exceptional education: at age seventeen, he was sent to Athens to study with Plato. He remained in Athens as a member of Plato's Academy for twenty years, until Plato's death about 347. Aristotle then spent several years in travel and study, crossing the Aegean Sea to Asia Minor and its coastal islands. During this period he undertook biological studies, and he encountered Theophrastus (from the island of Lesbos), who was to become his pupil and lifetime colleague. He returned to Macedonia in 342 to become the tutor of the young Alexander (later ``the Great''). In 335, when Athens fell under Macedonian rule, Aristotle returned to the city and began to teach in the Lyceum, a public garden frequented by teachers. He remained there, establishing an informal school, until shortly before his death in 322.

\zh 亚里士多德（图3.1）于公元前384年出生在希腊北部城镇斯塔吉拉，出身于一个有特权的家庭。他的父亲是马其顿国王阿明塔斯二世（亚历山大大帝的祖父）的私人医生。亚里士多德享有卓越教育的优势：十七岁时，他被送到雅典跟随柏拉图学习。他作为柏拉图学园的一员在雅典停留了二十年，直到柏拉图约于公元前347年去世。此后，亚里士多德用了数年时间旅行和研究，渡过爱琴海前往小亚细亚及其沿岸岛屿。在这一时期，他从事生物学研究，并遇到了泰奥弗拉斯托斯（来自莱斯博斯岛），后者后来成为他的学生和终身同事。公元前342年，他回到马其顿，担任年轻的亚历山大（后来被称为“大帝”）的导师。公元前335年，雅典处于马其顿统治之下时，亚里士多德回到这座城市，并开始在吕克昂（Lyceum）授课；吕克昂是一处教师们常去的公共园地。他在那里建立了一所非正式学校，并一直留到公元前322年去世前不久。

\orig In the course of his long career as student and teacher, Aristotle systematically and comprehensively addressed the major philosophical issues of his day. He is credited with more than 150 treatises, approximately 30 of which have come down to us. The surviving works appear to consist mainly of lecture notes or unfinished treatises not intended for wide circulation; whatever their exact origin, they were obviously directed to other philosophers, including advanced students. In modern translation, they occupy well over a foot of bookshelf, and they present a philosophical system overwhelming in power and scope. It is out of the question for us to survey the whole of Aristotle's philosophy, and we must be content with examining the fundamentals of his philosophy of nature---beginning with his response to positions taken by the pre-Socratics and Plato.

\zh 在其漫长的求学和教学经历中，亚里士多德系统而全面地处理了当时的主要哲学问题。传世目录归于他名下的论著超过150部，其中约30部流传至今。现存著作主要看来是讲义或未完成的论著，本不打算广泛流通；无论其确切来源如何，它们显然是写给其他哲学家看的，其中也包括高阶学生。按现代译本计算，它们占据书架的长度远不止一英尺，并呈现出一个力量和范围都令人震撼的哲学体系。我们不可能综览亚里士多德哲学的全貌，只能满足于考察其自然哲学的基本方面，先从他对前苏格拉底哲学家和柏拉图立场的回应开始。

\begin{figure}[htbp]
\centering
\includegraphics[width=0.42\linewidth]{../chunks/images/88ca2aa3a48f069b56deb949a5f0e30c086ed3d94423489f429999f388b54bf7.jpg}
\caption{Fig. 3.1. Aristotle. Museo Nazionale, Rome. Alinari/Art Resource N.Y. / 图3.1 亚里士多德。罗马，国立博物馆。Alinari/Art Resource N.Y.}
\end{figure}

\subsection*{Metaphysics and Epistemology / 形而上学与认识论}

\orig Through his long association with Plato, Aristotle had, of course, become thoroughly versed in Plato's theory of forms. Plato had drastically diminished (without totally rejecting) the reality of the material world observed by the senses. Reality in its perfect fullness, Plato argued, is found only in the eternal forms, which are dependent on nothing else for their existence. The objects that make up the sensible world, by contrast, derive their characteristics and their very being from the forms; it follows that sensible objects exist only derivatively or dependently.

\zh 由于长期与柏拉图相处，亚里士多德当然已经非常熟悉柏拉图的形式论。柏拉图大幅削弱了（但并未完全否定）感官所观察到的质料世界的实在性。柏拉图主张，充分而完美的实在只存在于永恒形式之中，而这些形式的存在不依赖任何其他东西。相比之下，构成可感世界的对象，其特征乃至其存在本身都来自形式；由此可见，可感对象只能以派生的或依赖的方式存在。

\orig Aristotle refused to accept this diminished, dependent status that Plato assigned to sensible objects. They must exist fully and independently, for in Aristotle's view they were what make up the real world. Moreover, the traits that give an individual object its character do not, Aristotle argued, have a prior and separate existence in a world of forms, but belong to the object itself. There is no perfect form of a dog, for example, existing independently in the world of forms and replicated imperfectly in individual dogs, imparting to them their attributes. For Aristotle, there were just individual dogs. These dogs certainly shared a set of attributes---for otherwise we would not be entitled to call them ``dogs''---but these attributes exist in, and belong to, individual dogs.

\zh 亚里士多德拒绝接受柏拉图赋予可感对象的这种被削弱的、依赖性的地位。可感对象必须充分且独立地存在，因为在亚里士多德看来，它们正是构成真实世界的东西。此外，亚里士多德论证说，赋予个别对象以特征的那些性状，并不在形式世界中具有先在而分离的存在，而是属于对象本身。比如，并不存在某个狗的完美形式，它独立存在于形式世界中，然后在个别的狗身上被不完美地复制，并把属性赋予它们。对亚里士多德来说，存在的只是个别的狗。这些狗当然共享一组属性，否则我们便无权称它们为“狗”；但这些属性存在于、并属于个别的狗。

\orig Perhaps this way of viewing the world has a familiar ring. Making individual sensible objects the primary realities (``substances,'' Aristotle called them) will seem like good common sense to most readers of this book, and probably struck Aristotle's contemporaries the same way. But if it makes good common sense, can it also be good philosophy? That is, can it deal successfully, or at least plausibly, with the difficult philosophical issues raised by the pre-Socratics and Plato---the nature of the fundamental reality, epistemological concerns, and the problem of change and stability? Let us take up these problems one by one.

\zh 这种看待世界的方式也许听起来很熟悉。把个别可感对象视为首要实在（亚里士多德称其为“实体”），对本书大多数读者来说会显得很合乎常识；对亚里士多德的同时代人来说，大概也是如此。但是，合乎常识是否也能成为好的哲学？也就是说，它能否成功地、或至少看似合理地处理前苏格拉底哲学家和柏拉图提出的那些困难哲学问题：根本实在的性质、认识论上的关切，以及变化与稳定的问题？让我们逐一讨论。

\orig The decision to locate reality in sensible, corporeal objects does not yet tell us very much about reality---only that we should look for it in the sensible world. Already in Aristotle's day, any philosopher would demand to know more: one thing he would demand to know was whether the corporeal materials of daily experience (wood, water, air, stone, metal, flesh, etc.) are themselves the fundamental, irreducible constituents of things, or whether they are composites of still more fundamental stuff. Aristotle addressed this question by drawing a distinction between properties and their subjects. He maintained (as most of us would) that a property has to be the property of something; we call that something its ``subject.'' To be a property is to belong to a subject; properties cannot exist independently.

\zh 决定把实在安置在可感的、有形的对象中，并没有告诉我们太多关于实在的内容；它只告诉我们应当在可感世界中寻找实在。早在亚里士多德时代，任何哲学家都会要求知道更多：其中一个问题便是，日常经验中的有形材料（木、水、气、石、金属、肉体等）本身是否就是事物的基本且不可还原的构成成分，还是说它们也是由更基本的东西复合而成。亚里士多德通过区分属性及其主体来处理这个问题。他主张（我们大多数人也会这样认为），属性必须是某物的属性；我们把这个“某物”称为属性的“主体”。作为属性就是属于某个主体；属性不能独立存在。

\orig Individual corporeal objects, then, have both properties (color, weight, texture, and the like) and something other than properties to serve as their subject. These two roles are played by ``form'' and ``matter,'' respectively. Corporeal objects are ``composites'' of form and matter---form consisting of the properties that make the thing what it is, matter serving as the subject or substratum for the form. A white rock, for example, is white, hard, heavy, and so forth, by virtue of its form; but matter must also be present, to serve as subject for the form, and this matter brings no properties of its own to its union with form. (Aristotle's doctrine will be further discussed in chap. 12, below, in connection with medieval attempts to clarify and extend it.)

\zh 因此，个别有形对象既具有属性（颜色、重量、质地等），也具有某种不是属性、但可以作为属性主体的东西。这两种角色分别由“形式”和“质料”承担。有形对象是形式与质料的“复合物”：形式由使事物成为其所是的那些属性构成，质料则作为形式的主体或基底。例如，一块白色的石头之所以是白的、硬的、重的等等，是凭借它的形式；但质料也必须在场，作为形式的主体，而这种质料在与形式结合时并不带来任何它自身的属性。（亚里士多德这一学说将在下文第12章中结合中世纪对它的澄清和扩展再作讨论。）

\orig We can never, in actuality, separate form and matter; they are presented to us only as a unitary composite. If they were separable, we should be able to put the properties (no longer the properties of anything) in one pile, the matter (absolutely propertyless) in another---an obvious impossibility. But if form and matter can never be separated, is it not meaningless to speak of them as the real constituents of things? Isn't this a purely logical distinction, existing in our minds, but not in the external world? Surely not for Aristotle, and perhaps not for us; most of us would think twice before denying the real existence of cold or red, although we can never collect a bucket of either one. In short, Aristotle once again surprises us by using commonsense notions to build a persuasive philosophical edifice.

\zh 实际上，我们永远不能把形式和质料分离开来；它们只作为一个统一复合物呈现在我们面前。如果二者可以分离，我们就应当能够把属性（不再是任何东西的属性）堆成一堆，把质料（绝对没有属性）堆成另一堆；这显然不可能。但是，如果形式和质料永远不能分离，那么把它们称为事物的真实构成要素，难道不是毫无意义吗？这难道不是一种纯粹逻辑上的区分，只存在于我们的心中，而不存在于外部世界吗？至少对亚里士多德来说绝非如此，也许对我们来说也不是；我们大多数人在否认“冷”或“红”的真实存在之前，都会三思，尽管我们永远不能收集一桶“冷”或一桶“红”。简言之，亚里士多德再次让我们惊讶：他用常识性的观念建构出一座有说服力的哲学大厦。

\orig Aristotle's claim that the primary realities are concrete individuals surely has epistemological implications, since true knowledge must be knowledge of truly real things. By this criterion, Plato's attention was naturally directed toward the eternal forms, knowable through reason or philosophical reflection. Aristotle's metaphysics of concrete individuals, by contrast, directed his quest for knowledge toward the material world of individuals, of nature, and of change---a world encountered through the senses.

\zh 亚里士多德关于首要实在是具体个体的主张，当然具有认识论含义，因为真正知识必须是关于真正实在之物的知识。按照这一标准，柏拉图的注意力自然转向通过理性或哲学反思而可知的永恒形式。相比之下，亚里士多德关于具体个体的形而上学，则把他对知识的追求引向由个体、自然和变化构成的质料世界，也就是通过感官遭遇到的世界。

\orig Aristotle's epistemology is complex and sophisticated. It must suffice here to indicate that the process of acquiring knowledge begins with sense experience. From repeated sense experience follows memory; and from memory, by a process of ``intuition'' or insight, the experienced investigator is able to discern the universal features of things. By the repeated observation of dogs, for example, an experienced dog breeder comes to know what a dog really is; that is, he comes to understand the form or definition of a dog, the crucial traits without which an animal cannot be a dog. Note that Aristotle, no less than Plato, was determined to grasp the universal traits or properties of things; but, unlike his teacher, Aristotle argued that one must start with the individual material thing. Once we grasp the universal properties or definition, we can put it to use as the premise of deductive demonstrations.

\zh 亚里士多德的认识论复杂而精细。这里只需指出，获得知识的过程始于感官经验。反复的感官经验产生记忆；由记忆出发，经过一种“直观”或洞见的过程，有经验的研究者便能辨识事物的普遍特征。举例说，通过反复观察狗，一位有经验的养犬者会认识到狗究竟是什么；也就是说，他会理解狗的形式或定义，即那些缺少了便不能称为狗的关键特征。需要注意，亚里士多德和柏拉图一样，决心把握事物的普遍性状或属性；但是，与其老师不同，亚里士多德主张必须从个别的质料事物出发。一旦我们把握了普遍属性或定义，就可以把它作为演绎证明的前提来使用。

\orig Knowledge is thus gained by a process that begins with experience (a term broad enough, in some contexts, to include common opinion or the reports of distant observers). In that sense knowledge is empirical; nothing can be known apart from such experience. But what we learn by this ``inductive'' process does not acquire the status of true knowledge until put into deductive form; the end product is a deductive demonstration (nicely illustrated in a Euclidean proof) beginning from universal definitions as premises. Although Aristotle discussed both the inductive and deductive phases (the latter far more than the former) in the acquisition of knowledge, he stopped considerably short of later methodologists, especially in the analysis of induction.

\zh 因此，知识是通过一个始于经验的过程获得的（“经验”一词在某些语境下足够宽泛，可以包括普通意见或远方观察者的报告）。在这个意义上，知识是经验性的；离开这种经验，什么也无法被认识。但是，我们通过这种“归纳”过程学到的东西，只有在被置入演绎形式之后，才获得真正知识的地位；其最终产物是从普遍定义作为前提出发的演绎证明（欧几里得式证明很好地说明了这一点）。虽然亚里士多德讨论了获得知识过程中的归纳阶段和演绎阶段（后者远多于前者），但他距离后来的方法论者仍相当遥远，尤其是在对归纳的分析方面。

\orig This is the theory of knowledge outlined by Aristotle in the abstract. Is it also the method actually employed in Aristotle's own scientific investigations? Probably not---with perhaps an occasional exception. Like modern scientists, Aristotle did not proceed by following a methodological recipe book, but rather by rough and ready methods, familiar procedures that had proved themselves in practice. Somebody has defined science as ``doing your damnedest, no holds barred''; when it came (for example) to his extensive biological researches, this is exactly what Aristotle did. It is not a surprise, and certainly no character defect, that Aristotle should, in the course of thinking about the nature and the foundations of knowledge, formulate a theoretical scheme (an epistemology) not perfectly consistent with his own scientific practice.

\zh 这就是亚里士多德在抽象层面勾勒出的知识理论。那么，这也是亚里士多德本人在科学研究中实际采用的方法吗？大概不是，或许偶有例外。像现代科学家一样，亚里士多德并不是照着一本方法论食谱做事，而是采用粗略但可行的方法，以及那些在实践中证明有效的熟悉程序。有人曾把科学定义为“竭尽全力，不设禁区”；就他广泛的生物学研究而言，亚里士多德正是这样做的。亚里士多德在思考知识的性质和基础时，提出一个与其自身科学实践并不完全一致的理论方案（认识论），这并不令人惊讶，当然也不构成性格缺陷。

\subsection*{Nature and Change / 自然与变化}

\orig The problem of change had become a celebrated philosophical issue (within the quite small community of philosophers) in the fifth century B.C. In the fourth century, Plato had dealt with it by restricting change to the imperfect material replica of the changeless world of forms. For Aristotle, a distinguished naturalist who was philosophically committed to the full reality of the changeable individuals that make up the sensible world, the problem of change was a most pressing one.

\zh 到公元前五世纪，变化问题已经成为一个著名的哲学议题（当然是在人数很少的哲学家共同体内部）。到公元前四世纪，柏拉图通过把变化限制在不变形式世界的不完美质料摹本中来处理这一问题。对亚里士多德这位杰出的自然研究者来说，变化问题尤其迫切，因为他在哲学上承认可感世界中那些变化着的个体具有充分实在性。

\orig Aristotle's starting point was the commonsense assumption that change is genuine. But this does not, by itself, get us very far; it remains to be demonstrated that the idea of change can withstand philosophical scrutiny; it must also be shown how change can be explained. Aristotle had various weapons in his arsenal by which to achieve these ends. The first was his doctrine of form and matter. If every object is constituted of form and matter, then Aristotle could make room for both change and stability by arguing that when an object undergoes change, its form changes (by a process of replacement, the new form replacing the old one) while its matter remains unchanged. Aristotle went on to argue that change in form takes place between a pair of opposites or contraries, one of which is the form to be achieved, the other its privation or absence. When the dry becomes wet, or the cold becomes hot, this is change from privation (dry or cold) to the intended form (wet or hot). Change, for Aristotle, is thus never random, but confined to the narrow corridor connecting pairs of contrary qualities; order is thus discernible even in the midst of change.

\zh 亚里士多德的出发点，是一个常识性假设：变化是真实的。但仅凭这一点并不能走得很远；还必须证明变化观念经得起哲学审查，也必须说明变化如何能够得到解释。为达到这些目的，亚里士多德有多种理论武器。第一种便是他的形式与质料学说。如果每个对象都由形式和质料构成，那么亚里士多德就可以通过如下论证为变化和稳定同时留出位置：当对象发生变化时，它的形式改变了（通过替换过程，新形式取代旧形式），而它的质料保持不变。亚里士多德进一步论证说，形式的变化发生在一对相反者或对立性质之间，其中一个是将要实现的形式，另一个则是它的缺失或不在。当干变成湿，或冷变成热时，这就是从缺失（干或冷）到预期形式（湿或热）的变化。因此，在亚里士多德看来，变化从不是随机的，而是被限制在连接成对对立性质的狭窄通道之中；即便在变化之中，秩序也清晰可辨。

\orig A determined Parmenidean might protest that to this point the analysis does nothing to escape Parmenides' objection to all change on the ground that inevitably it calls for the emergence of something out of nothing. Aristotle's reply is found in his doctrine of potentiality and actuality. Aristotle would undoubtedly have granted that if the only two possibilities are being and nonbeing---that is, if things either exist or do not exist---then the transition from non-hot to hot would indeed involve passage from nonbeing to being (the nonexistence of hot to the existence of hot) and would thus be vulnerable to Parmenides' objection. But Aristotle believed that the objection could be successfully circumvented by supposing that there are three categories associated with being instead of two: not just being and nonbeing, but (1) nonbeing, (2) potential being, and (3) actual being. If such is the state of things, then change can occur between potential being and actual being without nonbeing ever entering the picture. What Aristotle has in mind is perhaps most easily illustrated by examples from the biological realm. An acorn is potentially, but not actually, an oak tree. In becoming an oak tree, it becomes actually what it originally was only potentially. The change thus involves passage from potentiality to actuality---not from nonbeing to being, but from one kind or degree of being to another. Or for a pair of nonbiological examples, a heavy body held above the earth falls in order to fulfill its potential of being situated with other heavy things about the center of the universe. And a sculptor, with mallet and hammer, reveals in actuality a shape that existed potentially within the original block of marble.

\zh 一个坚定的巴门尼德信徒可能会抗议说，到目前为止，这一分析并没有避开巴门尼德对一切变化的反对：变化不可避免地要求某物从无中产生。亚里士多德的回答体现在他的潜能与现实学说中。亚里士多德无疑会承认，如果只有两种可能，即存在和非存在，或者说事物要么存在、要么不存在，那么从“非热”到“热”的转变确实会涉及从非存在到存在的过渡（从热的不存在到热的存在），因而会受到巴门尼德反驳的攻击。但是，亚里士多德认为，只要假定与存在相关的范畴不是两个而是三个，即不只是存在和非存在，而是（1）非存在，（2）潜在存在，（3）现实存在，就可以成功绕开这一反驳。若事情如此，变化就可以发生在潜在存在与现实存在之间，而无需让非存在进入图景。亚里士多德的意思，也许最容易用生物领域的例子来说明。一颗橡子潜在地是橡树，但现实地还不是橡树。在成为橡树时，它现实地成为了它原先只是潜在地所是的东西。因此，这种变化涉及从潜能到现实的过渡，不是从非存在到存在，而是从一种或一个程度的存在到另一种或另一个程度的存在。再举两个非生物的例子：悬在地面上方的重物下落，是为了实现它与其他重物一起位于宇宙中心附近的潜能；雕刻家用槌和凿，把原本潜在存在于大理石块中的形状现实地显现出来。

\orig If these arguments allow us to escape the logical dilemmas associated with the idea of change, and therefore to believe in its possibility, they do not yet tell us anything about the cause of change. Why should an acorn move from the status of potential oak tree to that of actual oak tree, or an object change from black to white, rather than remaining in its original state? Aristotle answered with an intricate, subtle, and not always consistent, theory of nature and causation. Given these difficulties, we will spare ourselves the pain of an exhaustive account and treat ourselves to the short version.

\zh 如果这些论证使我们得以摆脱变化观念所伴随的逻辑困境，从而相信变化是可能的，它们却尚未告诉我们变化的原因是什么。为什么橡子会从潜在橡树的状态转向现实橡树的状态？为什么一个对象会由黑变白，而不是停留在原有状态？亚里士多德用一套复杂、细致、且并非总是一贯的自然与因果理论来回答。鉴于这些困难，我们不必忍受详尽说明的痛苦，采用简短版本即可。

\orig The world we inhabit is an orderly one, in which things generally behave in predictable ways, Aristotle argued, because every natural object has a ``nature''---an attribute (associated primarily with form) that makes the object behave in its customary fashion, provided no insurmountable obstacle intervenes---or, as a modern commentator has put it, ``that within a thing which determines basically what that thing does when it is being itself.'' For Aristotle, a brilliant zoologist, the growth and development of biological organisms were easily explained by the activity of such an inner driving force. An acorn becomes an oak tree because its nature is to do so. But the theory was applicable beyond biological growth and, indeed, beyond the biological realm altogether. Dogs bark, rocks fall, and marble yields to the hammer and chisel of the sculptor because of their respective natures. Ultimately, Aristotle argued, all change and motion in the universe can be traced back to the natures of things. For the natural philosopher, who by definition is interested in change and things capable of undergoing change, these natures are the central object of study.

\zh 亚里士多德主张，我们所居住的世界是一个有秩序的世界，其中事物通常以可预期的方式行动，因为每个自然对象都有一种“自然本性”：这是一种属性（主要与形式相关），在没有不可克服的障碍介入时，使对象以其惯常方式行动；或者用一位现代评论者的话说，它是“事物内部那种基本决定该事物在成为其自身时会做什么的东西”。对杰出的动物学家亚里士多德来说，生物有机体的生长和发展，很容易由这种内在驱动力的活动来解释。橡子之所以成为橡树，是因为它的自然本性就是如此。但这一理论不仅适用于生物生长，事实上也不仅适用于生物领域。狗会吠叫，石头会下落，大理石会屈从于雕刻家的锤凿，都是由于它们各自的自然本性。亚里士多德最终主张，宇宙中所有变化和运动都可以追溯到事物的自然本性。对于自然哲学家而言，这些本性就是研究的中心对象；因为自然哲学家按定义所关心的，正是变化以及能够发生变化的事物。

\orig To this general statement of Aristotle's theory of ``nature,'' we need to add a qualification---namely, that an artificially produced object is a special case, for such an object possesses no nature other than the natures of its ingredients. If a chariot is constructed of wood and iron, the nature of wood and nature of iron do not yield to a composite ``nature of a chariot.'' By contrast, in the organic world the natures of the organs and tissues that make up an organism yield to the nature of the organism as a whole. The nature of the human body is not the sum of the natures of its various tissues and organs, but a unique nature characteristic of that living human as an organic whole.

\zh 对亚里士多德“自然”理论的这一概括性说明，还需要加上一项限定：人工制成的对象是特殊情况，因为这种对象除了其组成成分的自然本性之外，并不具有别的自然本性。如果一辆战车由木和铁构成，木的本性和铁的本性并不会让位于某种复合的“战车的本性”。相比之下，在有机世界中，组成有机体的器官和组织的本性，会让位于作为整体的有机体的本性。人体的本性不是其各种组织和器官本性的总和，而是作为有机整体的活生生的人所特有的一种独特本性。

\orig With this theory of nature in mind, we can understand a feature of Aristotle's scientific practice that has puzzled and distressed modern commentators and critics---namely, the absence from his work of anything resembling controlled experimentation. Unfortunately, such criticism overlooks Aristotle's aims, which drastically limited his methodological options. If, as Aristotle believed, the nature of a thing is to be discovered through the behavior of that thing in its natural, unfettered state, then artificial constraints will merely interfere and corrupt. If, despite interference, the object behaves in its customary fashion, we have troubled ourselves for no purpose. If we set up conditions that prevent the nature of an object from revealing itself, all we have learned is that it can be interfered with to the point of remaining concealed. Contrived experimentation violates, rather than reveals, the natures of things. Aristotle's scientific practice is not to be explained, therefore, as a result of stupidity or deficiency on his part---failure to perceive an obvious procedural improvement---but as a method compatible with the world as he perceived it and suited to the questions that interested him. Experimental science emerged not when, at long last, the human race produced somebody clever enough to perceive that artificial conditions would assist in the exploration of nature, but when a rich variety of conditions were fulfilled---including the emergence of questions to which such a procedure promised to provide answers.

\zh 牢记这一自然理论，我们就能理解亚里士多德科学实践中一个令现代评论者和批评者困惑并不满的特征：他的著作中缺少任何类似受控实验的东西。遗憾的是，这类批评忽视了亚里士多德的目标，而这些目标极大地限制了他可采用的方法。倘若如亚里士多德所信，事物的自然本性应当通过该事物在自然、无束缚状态下的行为来发现，那么人为约束只会干扰并破坏它。如果在干扰之下，对象仍按其惯常方式行动，我们便白费了力气；如果我们设置的条件阻止对象的自然本性显现，那么我们学到的也只是它可以被干扰到保持隐藏的程度。人为设计的实验不是揭示事物的自然本性，而是违背它们。因此，亚里士多德的科学实践不能解释为他的愚蠢或缺陷，也不是他未能看出某个显而易见的程序改进；它应被理解为一种与他所感知的世界相容、并适合他所关心问题的方法。实验科学之所以出现，并不是因为人类终于产生了一个聪明到足以看出人工条件有助于探索自然的人，而是因为一系列丰富条件得到满足，其中包括出现了这样一些问题：这种程序有望为它们提供答案。

\orig To complete our analysis of Aristotle's theory of change, we must briefly consider the celebrated four Aristotelian causes. To understand a change or the production of an artifact is to know its causes (perhaps best translated ``explanatory conditions and factors''). There are four of these: the form of a thing; the matter underlying that form, which persists through the change; the agency that brings about the change; and the purpose served by the change. These are called, respectively, formal cause, material cause, efficient cause, and final cause. To take an extremely simple example---the production of a statue---the formal cause is the shape given the marble, the material cause is the marble that receives this shape, the efficient cause is the sculptor, and the final cause is the purpose for which the statue is produced (perhaps the beautification of Athens or the celebration of one of its heroes). There are cases in which identifying one or another of the causes is difficult, or in which one or more causes merge, but Aristotle was convinced that his four causes provided an analytical scheme of general applicability.

\zh 为完成我们对亚里士多德变化理论的分析，必须简要考察著名的亚里士多德四因。理解一种变化或一件人工制品的产生，就是认识它的原因（也许最好译作“解释性条件和因素”）。这些原因有四种：事物的形式；承载该形式并在变化中持续存在的质料；引起变化的作用者；以及变化所服务的目的。它们分别称为形式因、质料因、动力因和目的因。举一个极为简单的例子：制作一尊雕像。形式因是赋予大理石的形状；质料因是接受这一形状的大理石；动力因是雕刻家；目的因则是制作这尊雕像的目的（也许是美化雅典，或纪念它的一位英雄）。在某些情况下，辨认某一原因会很困难，或者一个或多个原因会相互合并；但亚里士多德坚信，他的四因提供了一种具有普遍适用性的分析框架。

\orig We have said enough about the form-matter distinction to make clear what was meant by ``formal'' and ``material'' causes, and ``efficient'' cause is close enough to modern notions of causation to require no further comment; but ``final'' cause requires explanation. In the first place, the expression ``final cause'' is an English cognate derived from the Latin word \textit{finis}, meaning ``goal,'' ``purpose,'' or ``end,'' and it has nothing to do with the fact that it often appears last in the list of Aristotelian causes. Aristotle argued, quite rightly, that many things cannot be understood without knowledge of purpose or function. To explain the arrangement of teeth in the mouth, for example, we must understand their functions (sharp teeth in front for tearing, molars in back for grinding). Or to take an example from the inorganic realm, it is not possible to grasp why a saw is made as it is without knowing the function the saw is meant to serve. Aristotle went so far as to give final cause priority over material cause, noting that the purpose of the saw determines the material (iron) of which it must be made, whereas the fact that we possess a piece of iron does nothing to determine that we will make it into a saw.

\zh 关于形式与质料的区分，我们已经说得足够多，足以说明“形式因”和“质料因”的意思；“动力因”又与现代因果观念相当接近，无需进一步说明。但是，“目的因”需要解释。首先，“final cause”（目的因）这个英文表达源自拉丁词 \textit{finis}，意为“目标”“目的”或“终点”，它与目的因常常出现在亚里士多德诸因列表的最后这一事实没有关系。亚里士多德相当正确地论证说，许多事物若不了解其目的或功能，就无法理解。比如，要解释口腔中牙齿的排列，就必须理解它们的功能（前面的尖牙用于撕裂，后面的臼齿用于研磨）。再举一个无机领域的例子，如果不知道锯子要服务的功能，就不可能把握锯子为什么被制作成这个样子。亚里士多德甚至把目的因置于质料因之上，指出锯子的目的决定了它必须由什么材料（铁）制成；相反，我们拥有一块铁这一事实，并不能决定我们会把它做成锯子。

\orig Perhaps the most important point to be made about final cause is its clear illustration of the role of purpose (the more technical term is ``teleology'') in Aristotle's universe. The world of Aristotle is not the inert, mechanistic world of the atomists, in which the individual atom pursues its own course mindless of all others. Aristotle's world is not a world of chance and coincidence, but an orderly, organized world, a world of purpose, in which things develop toward ends determined by their natures. It would be unfair and pointless to judge Aristotle's success by the degree to which he anticipated modern science (as though his goal was to answer our questions, rather than his own); it is nonetheless worth noting that the emphasis on functional explanation to which Aristotle's teleology leads would prove to be of profound significance for all of the sciences and remains to this day a dominant mode of explanation within the biological sciences.

\zh 关于目的因，也许最重要的一点是：它清楚说明了目的（更技术性的术语是“目的论”）在亚里士多德宇宙中的作用。亚里士多德的世界并不是原子论者那种惰性的、机械论的世界；在原子论世界中，个别原子各行其道，对其他一切毫不关心。亚里士多德的世界不是偶然和巧合的世界，而是一个有秩序、有组织的世界，一个有目的的世界；在其中，事物依照其自然本性所规定的终点发展。以亚里士多德在多大程度上预示了现代科学来评判他的成功，是不公平且无意义的（仿佛他的目标是回答我们的问题，而不是他自己的问题）。尽管如此，仍值得注意的是，亚里士多德目的论所导向的功能性解释，对所有科学都将具有深远意义，并且至今仍是生物科学中的一种主导解释方式。

\subsection*{Cosmology / 宇宙论}

\orig Aristotle not only devised methods and principles by which to investigate and understand the world: form and matter, nature, potentiality and actuality, and the four causes. In the process, he also developed detailed and influential theories regarding an enormous range of natural phenomena, from the heavens above to the earth and its inhabitants below.

\zh 亚里士多德不仅提出了研究和理解世界的方法与原则：形式和质料、自然本性、潜能与现实，以及四因。在这一过程中，他还就范围极广的自然现象提出了详细而影响深远的理论，涵盖上方的天界以及下方的大地和其中的居民。

\orig Let us start with the question of origins. Aristotle adamantly denied the possibility of a beginning, insisting that the universe must be eternal. The alternative---that the universe came into being at some point in time---he regarded as unthinkable, violating (among other things) Parmenidean strictures about something coming from nothing. Aristotle's position on this question would prove troublesome for medieval Christian Aristotelians.

\zh 让我们从起源问题说起。亚里士多德坚决否认宇宙有开端的可能性，坚持认为宇宙必定是永恒的。另一种选择，即宇宙在时间中的某一点开始存在，被他视为不可思议；它违背了（除此之外还包括）巴门尼德关于某物不能从无中产生的禁令。亚里士多德在这一问题上的立场，后来会给中世纪基督教亚里士多德主义者带来麻烦。

\orig Aristotle considered this eternal universe to be a great sphere, divided into an upper and a lower region by the spherical shell in which the moon is situated. Above the moon is the celestial region; below is the terrestrial region; the moon, spatially intermediate, is also of intermediate nature. The terrestrial or sublunar region is characterized by birth, death, and transient change of all kinds; the celestial or supralunar region, by contrast, is a region of eternally unchanging cycles. That this scheme had its origin in observation would seem clear enough; in his \textit{On the Heavens}, Aristotle noted that ``in the whole range of time past, so far as our inherited records reach, no change appears to have taken place either in the whole scheme of the outermost heaven or in any of its proper parts.'' If in the heavens we observe eternally unvarying circular motion, he continued, we can infer that the heavens are not made of the terrestrial elements, the nature of which (observation reveals) is to rise or fall in transient rectilinear motions. The heavens must consist of an incorruptible fifth element (there are four terrestrial elements): the quintessence (literally, the fifth essence) or aether. The celestial region is completely filled with this quintessence (no void space) and divided, as we shall see, into concentric spherical shells bearing the planets. It had, for Aristotle, a superior, quasi-divine status.

\zh 亚里士多德认为，这个永恒宇宙是一个巨大的球体，并由月球所在的球壳分成上、下两个区域。月球之上是天界；月球之下是地界；月球在空间上居于中间，其性质也居于中间。地界或月下区域以生成、死亡以及各种短暂变化为特征；相比之下，天界或月上区域则是永恒不变循环的区域。这个图式源自观察，这一点似乎相当清楚；在《论天》中，亚里士多德指出：“在过去的整个时间范围内，就我们传承下来的记录所及，似乎无论在最外层天的整体结构中，还是在其任何固有部分中，都没有发生过变化。”他继续说，如果我们在天上观察到永恒不变的圆周运动，就可以推断天体不是由地上的元素构成的；因为观察显示，地上元素的本性是在短暂的直线运动中上升或下落。天体必定由一种不可朽坏的第五元素构成（地上元素有四种）：第五精华（字面意思即第五种本质），或以太。天界完全由这种第五元素充满（没有虚空），并且如我们将看到的，被划分为承载行星的同心球壳。对亚里士多德来说，天界具有一种较高的、近乎神圣的地位。

\orig The sublunar region is the scene of generation, corruption, and impermanence. Aristotle, like his predecessors, inquired into the basic element or elements to which the multitude of substances found in the terrestrial region can be reduced. He accepted the four elements originally proposed by Empedocles and subsequently adopted by Plato---earth, water, air, and fire. He agreed with Plato that these elements are in fact reducible to something even more fundamental; but he did not share Plato's mathematical inclination and therefore refused to accept Plato's regular solids and their constituent triangles. Instead, he expressed his own commitment to the reality of the world of sense experience by choosing sensible qualities as the ultimate building blocks. Two pairs of qualities are crucial: hot-cold and wet-dry. These combine in four pairs, each of which yields one of the elements (see fig. 3.2). Notice the use made once again of contraries. There is nothing to forbid any of the four qualities being replaced by its contrary, as the result of outside influence. If water is heated, so that the cold of water yields to hot, the water is transformed into air. Such a process easily explains changes of state (from solid to liquid to vapor, and conversely), but also more general transmutation of one substance into another. On such a theory as this, alchemists could easily build.

\zh 月下区域是生成、消灭和无常的场所。像他的前辈一样，亚里士多德探究了地界中众多物质可以被还原为何种基本元素。他接受了恩培多克勒最初提出、随后为柏拉图采纳的四元素：土、水、气、火。他同意柏拉图的看法，认为这些元素实际上还能还原为更基本的东西；但是，他并不分享柏拉图的数学倾向，因此拒绝接受柏拉图的正多面体及其构成三角形。相反，他选择可感性质作为最终构成要素，从而表达了自己对感官经验世界之实在性的承诺。两对性质至关重要：热--冷和湿--干。它们组合成四对，每一对产生一种元素（见图3.2）。请注意，对立性质在这里再次发挥作用。外部影响导致四种性质中的任何一种被其相反性质取代，这并无不可。若水被加热，使水的冷让位于热，水就转化为空气。这样的过程不仅容易解释状态变化（从固体到液体到气体，以及反向变化），也容易解释一种物质向另一种物质的更一般的嬗变。炼金术士很容易在这样一种理论上建立自己的学说。

\orig cold and dry = earth; cold and wet = water; hot and wet = air; hot and dry = fire.

\zh 冷而干 = 土；冷而湿 = 水；热而湿 = 气；热而干 = 火。

\begin{figure}[htbp]
\centering
\includegraphics[width=0.58\linewidth]{../chunks/images/749fd116e948c6da12cfa6db546ce82a7d13b303bbd503ba4bce425b8ec8e85a.jpg}
\caption{Fig. 3.2. Square of opposition of the Aristotelian elements and qualities. For a medieval (9th c.) version of this diagram, see John E. Murdoch, \textit{Album of Science: Antiquity and the Middle Ages}, p. 352. / 图3.2 亚里士多德元素与性质的对立方阵。关于此图的一个中世纪（9世纪）版本，见 John E. Murdoch, \textit{Album of Science: Antiquity and the Middle Ages}, p. 352。}
\end{figure}

\orig The various substances that make up the cosmos totally fill it, leaving no empty space. To appreciate Aristotle's view, we must lay aside our almost automatic inclination to think atomistically; we must conceive material things not as aggregates of tiny particles but as continuous wholes. If it is obvious that, say, a loaf of bread is composed of crumbs separated by small spaces, there is no reason not to suppose that those spaces are filled by some finer substance, such as air or water. And there is certainly no simple way of demonstrating, nor indeed any obvious reason for believing, that water and air are anything but continuous. Similar reasoning, applied to the whole of the universe, led Aristotle to the conclusion that the universe is full, a plenum, containing no void space. This claim would be attacked by medieval scholars.

\zh 构成宇宙的各种物质把宇宙完全充满，不留下任何空处。为了理解亚里士多德的观点，我们必须暂时放下那种几乎自动出现的原子式思维倾向；必须把质料事物设想为连续整体，而不是微小粒子的聚合体。比如，如果一条面包显然由许多面包屑组成，且面包屑之间隔着细小空隙，那么没有理由不设想这些空隙被某种更精细的物质，如空气或水，所填充。而且，也确实没有简单方法可以证明水和空气不是连续的；事实上，也没有明显理由相信它们不是连续的。类似推理应用到整个宇宙，便使亚里士多德得出结论：宇宙是充满的，是一个充盈体，没有虚空。这个主张后来会受到中世纪学者的攻击。

\orig Aristotle defended this conclusion with a variety of arguments, such as the following. The speed of a falling body is dependent on the density of the medium through which it falls---the less the density, the swifter the motion of the falling body. It follows that in a void space (density zero), there is nothing to slow the descent of the body, from which we would be forced to conclude that the body would fall with infinite speed---a nonsensical notion, since it implies that the body could be at two places at the same time. Critics have frequently noted that this argument can just as well be taken to prove that the absence of resistance does not entail infinite speed as to prove that void does not exist. The point is, of course, well taken. However, we need to understand that Aristotle's denial of the void did not rest on this single piece of reasoning. In fact, this was but one small part of a lengthy campaign against the atomists, in which Aristotle battled the notion of void space (or void place) with a variety of arguments, some more and some less persuasive.

\zh 亚里士多德用多种论证来捍卫这一结论，下面便是一例。落体的速度取决于它下落时穿过的介质密度；密度越小，落体运动越快。由此可推得，在虚空中（密度为零），没有任何东西减缓物体下落，于是我们就会被迫得出结论：物体将以无限速度下落。这个观念荒谬无稽，因为它意味着物体可以同时位于两个地方。批评者常常指出，这一论证同样可以被用来证明“没有阻力并不导致无限速度”，而不必证明虚空不存在。当然，这个批评切中要害。不过，我们需要理解，亚里士多德对虚空的否定并不依赖这一条推理。事实上，这只是他长期反对原子论者的一小部分；在这场论战中，亚里士多德用多种论证反击虚空空间（或虚空位置）的观念，其中有的较有说服力，有的较弱。

\orig In addition to being hot or cold and wet or dry, each of the elements is also heavy or light. Earth and water are heavy, but earth is the heavier of the two. Air and fire are light, fire being the lighter of the two. In assigning levity to two of the elements, Aristotle did not mean (as we might, if we were making the claim) simply that they are less heavy, but that they are light in an absolute sense; levity is not a weaker version of gravity, but its contrary. Because earth and water are heavy, it is their nature to descend toward the center of the universe; because air and fire are light, it is their nature to ascend toward the periphery (that is, the periphery of the terrestrial region, the spherical shell that contains the moon). If there were no hindrances, therefore, earth and water would collect at the center; because of its greater heaviness, earth would achieve a lower position, forming a sphere at the very center of the universe; water would collect in a concentric spherical shell just outside it. Air and fire naturally ascend, but fire, owing to its greater levity, occupies the outermost region, with air as a concentric sphere just inside it. In the ideal case (in which there are no mixed bodies and nothing prevents the natures of the four elements from fulfilling themselves), the elements would thus form a set of concentric spheres: fire on the outside, followed by air and water, and finally earth at the center (see fig. 3.3). But in reality, the world is composed largely of mixed bodies, one always interfering with another, and the ideal is never attained. Nonetheless, the ideal arrangement defines the natural place of each of the elements; the natural place of earth is at the center of the universe, of fire just inside the sphere of the moon, and so forth.

\zh 除了热或冷、湿或干之外，每一种元素也都是重的或轻的。土和水是重的，但土在二者中更重；气和火是轻的，火在二者中更轻。亚里士多德把轻性赋予两种元素时，并不是说（如我们若提出这一说法时可能会说的那样）它们只是较不重，而是说它们在绝对意义上是轻的；轻性不是重性的较弱版本，而是其相反者。由于土和水是重的，它们的自然本性便是下降到宇宙中心；由于气和火是轻的，它们的自然本性便是上升到外围（也就是地界的外围，包含月球的球壳）。因此，如果没有阻碍，土和水会聚集在中心；由于土更重，它会取得更低的位置，在宇宙最中心形成一个球体；水则会聚集在它外侧的同心球壳中。气和火自然上升；但火由于轻性更大，占据最外层区域，气则作为一个同心球层位于其内侧。在理想情况下（没有混合物体，也没有任何东西阻止四元素的本性实现自身），各元素便会形成一组同心球：外层是火，其次是气和水，最后是中心的土（见图3.3）。但在现实中，世界主要由混合物体构成，彼此总是相互干扰，理想状态从未实现。尽管如此，这种理想排列界定了每种元素的自然位置：土的自然位置在宇宙中心，火的自然位置就在月球天球内侧，等等。

\orig It must be emphasized that the arrangement of the elements is spherical. Earth collects at the center to form the earth, and it too is spherical. Aristotle defended this belief with a variety of arguments. Arguing from his natural philosophy, he pointed out that since the natural tendency of earth is to move toward the center of the universe, it must arrange itself symmetrically about that point. But he also called attention to observational evidence, including the circular shadow cast by the earth during a lunar eclipse and the fact that north-south motion by an observer on the surface of the earth alters the apparent position of the stars. Aristotle even reported an estimate by mathematicians of the earth's circumference (400,000 stades = about 45,000 miles, roughly 1.8 times the modern value). The sphericity of the earth, thus defended by Aristotle, would never be forgotten or seriously questioned. The widespread myth that medieval people believed in a flat earth is of modern origin.

\zh 必须强调，元素的排列是球形的。土聚集在中心形成大地，而大地本身也是球形的。亚里士多德用多种论证捍卫这一信念。从他的自然哲学出发，他指出，由于土的自然倾向是朝宇宙中心运动，它必定围绕这一点对称地排列自身。但他也提醒人们注意观察证据，包括月食时地球投下的圆形影子，以及地表观察者向南北方向移动会改变恒星的视位置这一事实。亚里士多德甚至报告了数学家对地球周长的一个估计（400,000斯塔德，约45,000英里，约为现代数值的1.8倍）。亚里士多德如此捍卫的地球球形说，后来从未被遗忘或受到严肃质疑。中世纪人相信地平说这一广泛流传的神话，乃是现代的产物。

\begin{figure}[htbp]
\centering
\includegraphics[width=0.58\linewidth]{../chunks/images/2800b314c989309a31b367a2dcd79da8bd2197dacae4678ce6e7f2442cd1e18d.jpg}
\caption{Fig. 3.3. The Aristotelian cosmos. / 图3.3 亚里士多德宇宙。}
\end{figure}

\orig Finally, we must note one of the implications of this cosmology, namely that space, instead of being a neutral, homogeneous backdrop (analogous to our modern notion of geometrical space) against which events occur, has properties. Or to express the point more precisely, ours is a world of space, whereas Aristotle's was a world of place. Heavy bodies move toward their place at the center of the universe not because of a tendency to unite with other heavy bodies located there, but simply because it is their nature to seek that central place; if by some miracle the center happened to be vacant (a physical impossibility in an Aristotelian universe, but an interesting imaginary state of affairs), it would remain the destination of every heavy body.

\zh 最后，我们必须注意这一宇宙论的一个含义：空间并不是事件发生时一个中性的、同质的背景（类似我们现代的几何空间观念），而是具有属性。或者更准确地说，我们的世界是空间的世界，而亚里士多德的世界则是位置的世界。重物朝向宇宙中心的自然位置运动，并不是因为它们倾向于与那里其他重物结合，而只是因为它们的自然本性就是寻求那个中心位置；即使奇迹般地中心恰好空着（这在亚里士多德宇宙中是物理上不可能的，但作为想象情形颇有趣），它仍会是每一个重物的目的地。

\subsection*{Motion, Terrestrial and Celestial / 地上运动与天体运动}

\orig We can best understand Aristotle's theory of motion by grasping its two most fundamental claims. The first is that motion is never spontaneous; there is no motion without a mover. The second is the distinction between two types of motion: motion toward the natural place of the moving body is ``natural'' motion; motion in any other direction occurs only under coercion from an outside force and is therefore a ``forced'' or ``violent'' motion.

\zh 要最好地理解亚里士多德的运动理论，需把握其两个最基本的主张。第一，运动从不是自发的；没有推动者，就没有运动。第二，运动分为两类：朝向运动物体自然位置的运动，是“自然”运动；任何其他方向的运动，都只能在外力强迫下发生，因此是“强迫的”或“暴力的”运动。

\orig The mover in the case of natural motion is the nature of the body, which is responsible for its tendency to move toward its natural place as defined by the ideal spherical arrangement of the elements. Mixed bodies have a directional tendency that depends on the proportion of the various elements in their composition. When a body undergoing natural motion reaches its natural place, its motion ceases. The mover in the case of forced motion is an external force, which compels the body to violate its natural tendency and move in a direction or manner other than straight-line motion toward its natural place. Such motion ceases when the external force is withdrawn.

\zh 在自然运动的情况下，推动者是物体的自然本性；这种本性负责使物体倾向于朝向它在元素理想球形排列中所规定的自然位置运动。混合物体的方向倾向取决于其构成中各种元素的比例。当一个进行自然运动的物体到达其自然位置时，它的运动便停止。在强迫运动的情况下，推动者是外在力量；它迫使物体违背其自然倾向，以不同于朝向自然位置的直线运动的方向或方式运动。当外力撤去时，这种运动便停止。

\orig So far, this seems sensible. One obvious difficulty, however, is to explain why a projectile hurled horizontally, and therefore undergoing forced motion, does not come to an immediate halt when it loses contact with whatever propelled it. Aristotle's answer was that the medium takes over as mover. When we project an object, we also act on the surrounding medium (air, for instance), imparting to it the power to move objects; this power is communicated from part to part, in such a way that the projectile is always in contact with a portion of the medium capable of keeping it in motion. If this seems implausible, consider the greater implausibility (from Aristotle's standpoint) of the alternative---that a projectile, which is inclined by nature to move toward the center of the universe, moves horizontally or upward despite the fact that there is no longer anything causing it to do so.

\zh 到目前为止，这似乎合情合理。然而，一个明显困难是解释为什么一个被水平抛出的投射物，既然处于强迫运动之中，在失去与推动它的东西的接触后并不会立即停下。亚里士多德的回答是，介质接替成为推动者。当我们投射一个物体时，也作用于周围介质（例如空气），把移动物体的能力赋予介质；这种能力从一个部分传递到另一个部分，使投射物总是与某部分能够维持其运动的介质相接触。如果这看起来难以置信，那么请考虑另一种选择从亚里士多德立场看更加难以置信：一个自然上倾向于朝宇宙中心运动的投射物，竟然在已经没有任何东西使它这样做的情况下，仍然水平或向上运动。

\orig Force is not the only determinant of motion. In all real cases of motion in the terrestrial realm, there will also be a resistance or opposing force. And it seemed clear to Aristotle that the quickness of motion must depend on these two determining factors---the motive force and the resistance. The question arose: what is the relationship between force, resistance, and speed? Although it probably did not occur to Aristotle that there might be a quantitative law of universal applicability, he was not without interest in the question and did make several forays into quantitative territory. In reference to natural motion in his \textit{On the Heavens} and again in his \textit{Physics}, Aristotle claimed that when two bodies of differing weight descend, the times required to cover a given distance will be inversely proportional to the weights. (A body twice as heavy will require half the time.) In the same chapter of the \textit{Physics}, Aristotle introduced resistance into the analysis of natural motion, arguing that if bodies of equal weight move through media of different densities, the times required to traverse a given distance are proportional to the densities of the respective media; that is, the greater the resistance the slower the body moves. Finally, Aristotle also dealt with forced motion in his \textit{Physics}, claiming that if a given force moves a given weight (against its nature) for a given distance in a given time, the same force will move half that weight twice the distance in that same time (or the same distance in half that time); alternatively, half the force will move half the weight the same distance in the same time.

\zh 力并不是运动的唯一决定因素。在地界中一切真实运动的情形里，也会存在阻力或反向力。在亚里士多德看来，运动快慢显然必须取决于这两个决定因素：推动力和阻力。问题随之出现：力、阻力和速度之间的关系是什么？虽然亚里士多德大概没有想到可能存在一个普遍适用的定量规律，但他并非对此问题毫无兴趣，并且确实多次涉足定量领域。关于自然运动，亚里士多德在《论天》中、并且又在《物理学》中主张，当两个重量不同的物体下落时，它们通过给定距离所需时间与其重量成反比。（一个重两倍的物体将需要一半时间。）在《物理学》的同一章中，亚里士多德把阻力引入自然运动分析，论证说，如果同等重量的物体穿过不同密度的介质运动，那么它们通过给定距离所需时间与相应介质的密度成正比；也就是说，阻力越大，物体运动越慢。最后，亚里士多德在《物理学》中也处理了强迫运动，声称如果某一给定力能在给定时间内使某一给定重量（违反其本性）移动给定距离，那么同样的力就能在同样时间内使一半重量移动两倍距离（或在一半时间内移动同样距离）；或者，半个力能使一半重量在同样时间内移动同样距离。

\orig From such statements, some of Aristotle's successors have made a determined effort to extract a general law. This law is customarily stated as:
\[
\mathrm{v} \propto \mathrm{F}/\mathrm{R}.
\]
That is, velocity (v) is proportional to the motive force (F) and inversely proportional to the resistance (R). For the special case of the natural descent of a heavy body, the motive force is the weight (W) of the body, and the relationship then becomes:
\[
\mathrm{v} \propto \mathrm{W}/\mathrm{R}.
\]

\zh 从这些陈述中，亚里士多德的一些后继者努力抽取出一个普遍规律。这个规律通常表述为：
\[
\mathrm{v} \propto \mathrm{F}/\mathrm{R}.
\]
也就是说，速度（v）与推动力（F）成正比，并与阻力（R）成反比。对于重物自然下落这一特殊情况，推动力就是物体的重量（W），于是关系变为：
\[
\mathrm{v} \propto \mathrm{W}/\mathrm{R}.
\]

\orig Such relationships probably do no great violence to Aristotle's intent for most cases of motion; however, giving them mathematical form, as we have done, suggests that they hold for all values of v, F (or W), and R---a claim that Aristotle would certainly have denied. He stated explicitly, for example, that a resistance equal to the motive force will prevent motion altogether, whereas the formula above offers no such result. Moreover, the appearance of velocity in these relationships seriously misrepresents Aristotle's conceptual framework, which contained no concept of velocity as a quantifiable measure of motion, but described motion only in terms of distances and times. Velocity as a technical scientific term to which numerical values might be assigned was a contribution of the Middle Ages (see below, chap. 12).

\zh 对于大多数运动情形来说，这些关系大概并没有严重违背亚里士多德的意图；然而，像我们这样赋予它们数学形式，便暗示这些关系对 v、F（或 W）和 R 的所有数值都成立，而亚里士多德肯定会否认这一点。例如，他明确说过，若阻力等于推动力，运动就会完全被阻止；而上述公式并不给出这种结果。此外，速度出现在这些关系中，也严重误表了亚里士多德的概念框架；他的框架中并没有把速度作为运动的可量化尺度这一概念，而只是用距离和时间来描述运动。速度作为可以赋予数值的技术性科学术语，是中世纪的贡献（见下文第12章）。

\orig Aristotle has been severely criticized for this theory of motion, on the assumption that any sensible person should have recognized its fatal flaws. Is such criticism justified? In the first place, our goal is to understand the behavior, beliefs, and achievements of historical actors against the background of the culture in which they lived, rather than to assess credit or blame according to the degree to which those historical actors resemble us. In short, historians must always contextualize their subjects. Second, some of the criticisms of Aristotle's theories of motion apply only to theories foisted onto Aristotle by followers and critics, rather than to his own. Third, the theory in its genuinely Aristotelian (and properly contextualized) version makes quite good sense today and would surely have made good sense in the fourth century B.C. For example, various surveys have shown that the majority of modern, university-educated people are prepared to assent to many of the basics of Aristotle's theory of motion. Fourth, the relatively modest level of quantitative content in Aristotle's theory is easily explained as the outcome of his larger philosophy of nature. His primary goal was to understand essential natures, not to explore quantitative relationships between such incidental factors as the space-time (or place-time) coordinates applicable to a moving body; even an exhaustive investigation of the latter gives us no useful information about the former. You may criticize Aristotle, if you like, for not being interested in whatever interests modern scientists, but we do not thereby learn anything significant about Aristotle.

\zh 亚里士多德因其运动理论受到严厉批评，批评者假定任何明智的人都应当看出其中致命缺陷。这种批评合理吗？首先，我们的目标是在历史行动者所生活的文化背景中理解他们的行为、信念和成就，而不是按照他们在多大程度上像我们来评判功过。简言之，历史学家必须始终把研究对象放回语境中。第二，对亚里士多德运动理论的一些批评，只适用于后来追随者和批评者强加给亚里士多德的理论，而不适用于他自己的理论。第三，在真正亚里士多德式的（并且经过恰当语境化的）版本中，这一理论今天也相当说得通，在公元前四世纪当然更会显得合情合理。例如，多项调查显示，现代受过大学教育的人中，大多数愿意赞同亚里士多德运动理论的许多基本点。第四，亚里士多德理论中定量内容相对有限，很容易解释为他更大范围自然哲学的结果。他的主要目标是理解本质性的自然本性，而不是探索运动物体所适用的时空（或位置--时间）坐标等偶然因素之间的定量关系；即使把后者研究得再详尽，也不会给我们关于前者的有用信息。如果你愿意，可以批评亚里士多德不关心现代科学家所关心的东西，但我们并不能因此学到任何关于亚里士多德的重要内容。

\orig Motion in the celestial sphere is an altogether different sort of phenomenon. The heavens, composed of the incorruptible quintessence, possess no contraries and are therefore incapable of qualitative change. It might seem fitting for such a region to be absolutely motionless, but this hypothesis is defeated by the most casual observation of the heavens. Aristotle therefore assigned to the heavens the most perfect of motions---continuous uniform circular motion. Besides being the most perfect of motions, uniform circular motion appears to have the capability of explaining the observed celestial cycles.

\zh 天球中的运动则是完全不同的一类现象。天界由不可朽坏的第五元素构成，没有相反者，因此不能发生性质变化。这样一个区域似乎应当绝对静止，但对天空最随意的观察都会推翻这一假设。因此，亚里士多德把最完美的运动赋予天界，即连续的匀速圆周运动。匀速圆周运动除了是最完美的运动之外，似乎还能够解释观察到的天体周期。

\orig By Aristotle's day, these cycles had been an object of study for centuries in the Greek world and for millennia in its predecessor civilizations. It was understood that the ``fixed'' stars move with perfect uniformity, as though fixed to a uniformly rotating sphere, with a period of rotation of approximately one day. But there were seven stars, the wandering stars or planets, that displayed a more intricate motion, apparently crawling around on the stellar sphere as it went through its daily rotation. These seven were the Sun, Moon, Mercury, Venus, Mars, Jupiter, and Saturn. The sun crawls slowly (about $1^\circ$/day), west to east with small variations in speed, through the sphere of fixed stars along a path called the ecliptic, which passes through the center of the zodiac. The moon follows approximately the same course, but at the more rapid rate of about $12^\circ$/day. The remaining planets also move along the ecliptic (or in its vicinity) with variable speed and with an occasional reversal of direction.

\zh 到亚里士多德时代，这些周期在希腊世界已经被研究了数百年，在希腊之前的文明中则已被研究了数千年。人们理解到，“恒定的”星辰以完美的均匀性运动，仿佛固定在一个匀速旋转的天球上，旋转周期约为一日。但是，有七颗星，即游星或行星，表现出更复杂的运动：它们显然在恒星天球进行每日旋转时，于其上缓慢爬行。这七颗星是太阳、月亮、水星、金星、火星、木星和土星。太阳沿着一条称为黄道、穿过黄道带中心的路径，在恒星天球上由西向东缓慢移动（约每天 $1^\circ$），速度有小幅变化。月亮大致沿同一路径运行，但速度更快，约每天 $12^\circ$。其余行星也沿黄道（或其附近）运动，速度可变，并且偶尔逆行。

\orig Are such complex motions compatible with the requirement of uniform circular motion in the heavens? Eudoxus, a generation before Aristotle, had already shown that they are. I will return to this subject in chap. 5; for the moment, it will be sufficient to point out that Eudoxus treated each complex planetary motion as a composite of a series of simple uniform circular movements. He did this by assigning to each planet a set of concentric spheres, and to each sphere one component of the complex planetary motion. Aristotle took over this scheme, with various modifications. When he was finished, he had produced an intricate piece of celestial machinery, consisting of fifty-five planetary spheres plus the sphere of the fixed stars.

\zh 这种复杂运动与天界中匀速圆周运动的要求相容吗？比亚里士多德早一代的欧多克索斯已经表明二者相容。第5章将回到这个主题；眼下只需指出，欧多克索斯把每一种复杂的行星运动处理为一系列简单匀速圆周运动的复合。他的做法是为每颗行星分配一组同心球，并把复杂行星运动的一个分量分配给每个球。亚里士多德接过这一方案，并作了若干修改。完成时，他已经构造出一套精巧的天体机械，由五十五个行星天球加上恒星天球构成。

\orig What is the cause of movement in the heavens? Aristotle's natural philosophy would not allow such a question to go unasked. The celestial spheres are composed, of course, of the quintessence; their motion, being eternal, must be natural rather than forced. The cause of this eternal motion must itself be unmoved, for if we do not postulate an unmoved mover, we quickly find ourselves trapped in an infinite regress: a moving mover must have acquired its motion from yet another moving mover, and so on. Aristotle identified the unmoved mover for the planetary spheres as the ``Prime Mover,'' a living deity representing the highest good, wholly actualized, totally absorbed in self-contemplation, nonspatial, separated from the spheres it (or he or she) moves, and not at all like the traditional anthropomorphic Greek gods. How, then, does the Prime Mover or Unmoved Mover cause motion in the heavens? Not as efficient cause, for that would require contact between the mover and the moved, but as final cause. That is, the Prime Mover is the object of desire for the celestial spheres, which endeavor to imitate its changeless perfection by assuming eternal, uniform circular motions. Any reader who has followed this much of Aristotle's discussion would be justified in assuming that there is a single Unmoved Mover for the entire cosmos; it comes as something of a surprise, therefore, when Aristotle announces that, in fact, each of the celestial spheres has its own Unmoved Mover, the object of its affection and the final cause of its motion.

\zh 天上的运动原因是什么？亚里士多德的自然哲学不会允许这个问题不被提出。天球当然由第五元素构成；它们的运动既然是永恒的，就必定是自然运动而不是强迫运动。这种永恒运动的原因本身必须是不动的，因为如果不设定一个不动的推动者，我们很快就会陷入无穷倒退：一个运动着的推动者必定从另一个运动着的推动者那里获得其运动，如此无穷。亚里士多德把行星天球的不动推动者认定为“第一推动者”：这是一位活的神祇，代表最高善，完全现实化，完全沉浸于自我观照，没有空间位置，与它（或他、或她）所推动的天球相分离，并且完全不像传统拟人化的希腊诸神。那么，第一推动者或不动的推动者如何造成天上的运动？不是作为动力因，因为那要求推动者与被推动者之间有接触；而是作为目的因。也就是说，第一推动者是天球欲求的对象；天球试图通过采取永恒、匀速的圆周运动来模仿其不变的完美。任何读到亚里士多德这一部分讨论的读者，都有理由假定整个宇宙只有一个不动的推动者；因此，当亚里士多德宣布事实上每一个天球都有自己的不动推动者，即它所钟爱的对象和其运动的目的因时，这多少令人惊讶。

\end{document}
